\section{Introduction}

A plane graph is a specific topological embedding of a planar graph in a two-dimensional Euclidean space, where vertices are mapped to distinct coordinates and edges are mapped to continuous arcs that intersect exclusively at their common endpoints. Identifying and extracting the discrete regions (or faces) bounded by these edges is a fundamental operation across an extensive array of domains, including geographic information systems (GIS), computer graphics, VLSI layout verification, mesh generation, and automated robotics navigation~\cite{ANTOLIN2023103448,DEMAGALHAES2020102801,10.1145/3139958.3140001}. In modern computational engineering, simulations involving finite element analysis or fluid dynamics frequently manipulate structural meshes containing tens of millions of distinct elements~\cite{Lintermann2021,LINTERMANN2014131,electronics14153151}. At this scale, serial or poorly parallelized geometric extraction routines quickly become severe computational bottlenecks.

% The first optimal sequential algorithm for region extraction that will bed discussed in the paper was proposed Mark de Berg et al~\cite{deberg2008computational} and it is based on Doubly Connected Edge List (DCEL) data structure and the second was proposed by Jiang and Bunke~\cite{JIANG1993553} and is based in Wedges data structure. To differentiate between them the algorithm will be called DCEL-based and Wedge-based algorithms, respectively. Both of them are achieving a time complexity of \(O(m \log m)\) and space complexity of \(O(m)\), where \(m\) is the number of edges in the graph \(G(V, E)\). As shown by Euler's formula for planar graphs, a simple plane graph with \(n \geq 3\) vertices has at most \(m \leq 3n - 6\) edges~\cite{Diestel2025}.

Two primary optimal sequential algorithms for region extraction are discussed in this paper. The first one~\cite{JIANG1993553} is based on wedges, while the second one~\cite{deberg2008computational} is based on half-edges (HEs) and is implemented in many computational geometry libraries, such as the Computational Geometry Algorithms Library (CGAL)~\cite{10.1145/1653771.1653865} and the Boost Graph Library (BGL)~\cite{10.5555/504206}. In this paper, these approaches are referred to as the Wedge-based and HE-based algorithms, respectively. Both methods achieve a time complexity of $O(m \log m)$ and a space complexity of $O(m)$, where $m$ denotes the number of edges in the graph $G(V, E)$. As shown by Euler's formula for planar graphs, a simple plane graph with $n \geq 3$ vertices has at most $m \leq 3n - 6$ edges \cite{Diestel2025}, confirming that such graphs are inherently sparse.

Although theoretically optimal, the original sequential algorithms can be computationally impractical when applied to large-scale graphs commonly encountered in modern scientific and engineering applications, and the emergence of massively parallel platforms, such as GPUs with NVIDIA's CUDA~\cite{10.1145/1401132.1401152} framework, creates new opportunities for accelerating algorithmic graph operations.

% In this paper, we propose a CUDA-based GPU implementation of the optimal region extraction algorithm for plane graphs, emphasizing the design of data structures and execution strategies optimized for GPU architectures. The GPU implementation is evaluated against different reference CPU versions based on the Computational Geometry Algorithms Library (CGAL)~\cite{10.1145/1653771.1653865}, the Boost Graph Library (BGL)~\cite{10.5555/504206}, and our own implementation of the algorithms, using a suite of synthetic 2D mesh datasets that simulate various planar graph structures. Experimental results confirm that the GPU implementation maintains correctness while achieving a substantial speedup compared to the CPU baseline. These findings highlight the effectiveness of parallelization for graph-based geometric processing tasks and demonstrate the practicality of the approach for large-scale applications.

This paper introduces comprehensive GPU-accelerated parallel adaptations of both the Wedge-based and the Half-Edge-based region extraction algorithms. We evaluate our parallel CUDA C++ implementations against established CPU geometric frameworks, including CGAL and BGL, as well as heavily optimized sequential baselines. Our parallel designs adapt a parallel pointer-jumping (path doubling)~\cite{jaja1992} technique and leverage the parallel primitives of the Thrust library~\cite{BELL2012359}. We utilize a suite of synthetic 2D mesh datasets to simulate the large-scale planar graph structures encountered in modern engineering. Experimental results demonstrate that our GPU implementations maintain absolute correctness while achieving substantial speedups—up to 33$\times$ in our tests—compared to CPU baselines. These findings highlight the effectiveness of parallelization for geometric graph processing and demonstrate the practicality of our approach for large-scale applications.

The remainder of this work is structured as follows. Section~\textsc{ii} formalizes the necessary mathematical definitions and structural notations. Section~\textsc{iii} outlines the operational logic of the baseline sequential algorithms. Section~\textsc{iv} provides a deep overview of our proposed parallelization methodologies and CUDA implementation strategies. Section~\textsc{v} presents the experimental results, performance characteristics, and workload profiling, and Section~\textsc{vi} concludes the paper.